\chaptertitle{8}{挑选指数基金的三个秘诀}

在前几章里，螺丝钉介绍了几十个指数，围绕这些指数又有不同的基金公司开发的几百只产品。

追踪一个指数的指数基金，数量可能很多，我们究竟要选择什么样的指数基金来投资呢？挑选指数基金有哪些需要注意的地方？

这一章，我们就来了解一下。

有的指数比较小众，可能对应的指数基金就只有一只。这种挑选起来会比较容易。但是像沪深300、中证500、上证50等指数，对应的指数基金品种很多。

如果我们搜索市面上的沪深300指数基金，会发现有几十只沪深300指数基金。几乎每家大型基金公司都有沪深300指数基金。

因为沪深300指数的规则是公开透明的，不同基金公司都可以围绕沪深300指数开发自己的产品。时间久了，市面上的沪深300指数基金数量就会非常多。这给我们的挑选带来一些困难。

挑选指数基金的时候，通常有三个比较重要的因素，分别是规模、费用和追踪误差。

\subsection{秘诀一：基金规模小于1亿元不要选}
{\heiti 什么是基金的规模}

基金是基金公司开发出来的产品，在基金刚成立的认购期，以及基金成立之后开放申购期，我们都可以买到基金。这些资金汇集起来，就到了基金对应的托管账户中，最后基金经理用这些钱去购买股票、债券等资产。

一只基金所持有的股票、债券、现金等的总市值，就是这只基金的规模。

通常来说，我们不考虑规模太小的基金。

如果一只指数基金规模较小，它清盘的概率就较大。基金清盘并不是说我们的投资血本无归了，而是按照某一基金净值强制赎回，导致我们的投资中断。如果基金规模太小，那么基金公司运作这只基金可能就是亏本的，基金公司就有可能停止这只基金的运作。所以一般挑选指数基金的时候，会避开规模较小的指数基金，最好规模在1亿元以上再考虑。

{\heiti 什么是基金的清盘}

基金公司是商业机构，也有运营成本的。如果基金公司不赚钱，就没有动力宣传，没有动力宣传基金规模就越小，导致恶性循环。

基金规模如果太小了，基金费率就会覆盖不了成本，那运营下去也就没有意思了。特别是一些比较冷门的指数基金，更是如此。所以，小基金有清盘的风险。

清盘不是说投资者血本无归。实际上清盘之前，基金公司会留出时间让投资者手动赎回要清盘的基金。即使开始走清盘流程了，最后基金公司也会按照一个净值强制赎回，把资金归还给投资者。

但这毕竟比较麻烦，影响我们的投资。所以通常我们不考虑规模小于1亿元的指数基金品种。

要注意的是，基金规模小于1亿元，并不意味着基金一定会马上开始清盘。基金公司也有可能愿意亏本运营基金。

截至2019年3月，国内有651只指数基金，但是有296只指数基金，规模小于1亿元。规模2亿元以上的，只有255只。

很多人都不太了解指数基金，导致很多质量不错的指数基金最后的规模很小。如果基金公司运作长期亏损，这只基金就逐渐清盘了。

{\heiti 基金规模大小对投资的影响}

基金规模小，还有一些缺点。

基金运作的时候，会有一些固定费用。这些费用的金额是固定的，有的是营销费用，有的是律师会计等的费用，这些不管是大基金还是小基金都会有。如果基金规模较小，这些固定费用在基金净值中的占比就会比较高，影响收益。

例如每年固定费用如果是50万元，对一只规模5000万元的基金来说，50万元就占了1\%，每年拖累基金1\%的收益。但如果基金规模是5亿元，那50万元就只占0.1\%，影响就小很多了。

基金规模小，还有可能在基金申购赎回的时候遇到一些流动性问题。例如规模太小的基金，如果遇到持有者大比例赎回，可能赎回到账时间就会比一般的赎回到账时间晚一些，虽然最后也是可以赎回的。

所以，我们挑选基金时，通常不考虑规模1亿元以下的基金品种。

那基金规模非常大呢？

如果指数基金的规模非常大，那理论上会对指数基金追踪指数产生麻烦。比如说一个指数背后的公司有100家，平均每家公司的规模是10亿元，这只指数基金，理论上的规模上限就是1000亿元，因为背后公司股票总市值就这么大。给基金再多的钱，也买不到这个指数的股票了。

不过现阶段，在国内还不用担心。国内的指数基金品种还不太普及，指数基金总规模很小。大部分指数基金的规模都是几亿十几亿元。所以很长时间里，规模过大的问题都不需要担心。

{\heiti 基金会不会“跑路、爆雷”}

这几年国内也出现过一些跑路、爆雷的金融产品，也有的朋友会有疑问，原来基金也会有清盘，那基金本身是不是足够安全呢？

首先我们要明白，什么是跑路，什么是爆雷。

什么是跑路呢？

就是我们投资了某个投资品种，然后这个品种卷钱消失了，导致投资者产生巨大的亏损。

什么是爆雷呢？

比如，P2P把钱贷了出去，结果钱收不回来，资金周转不过来，导致无法支付投资者本金利息。

指数基金是不会出现这些风险的。为什么呢？

第一，成立公募基金公司门槛高。

在国内，想要发行公募基金，需要有公募基金牌照。

成立一家公募基金公司，得有起码1亿元的注册资本，并且是实缴的。其实，现在就算有2亿元的资本，也不一定能成立一家公募基金公司。而且这仅仅是最基本的要求，还要求股东资产规模要达标，并且最近3年没有任何不良信用记录。

所以，成立一家公募基金公司门槛还是挺高的，而且成立公募基金公司，就已经受到了严格的监管。

第二，任何一只公募基金都有托管人。

也就是说，一只公募基金产品，资金并不是直接存放在基金公司，而是放在托管银行。

我们看到公募基金的托管费，其实就是托管银行收取的。并不是所有银行都有资格成为托管银行。国内规模较大的头部银行才有托管资格。另外，国内规模最大的几家券商也有托管资格。托管银行或者托管券商的数量，比国内公募基金公司的数量还少。

我们申购基金的资金，都是放在这个托管账户上。基金经理可以调动这些资金去买股票、债券，但是卖出股票、债券所得的资金，只能放在这个托管账户上。

即使是基金公司，也没有资格把这些资金转到其他账户。所以这种托管机制，决定了公募基金是不可能出现跑路的情况的。

第三，公募基金受到证监会、基金业协会的严格监管。

无论是资金的投资，或是风险控制，或是整个投资的流程，都有专门的部门和负责人负责监控。公募基金公司自己也会做好风险控制。

虽然基金投资可能会产生亏损，但是从产品本身的结构上，不会出现基金公司卷钱跑路的情况。

我们知道了，公募基金不会跑路，那么会不会出现爆雷呢？

P2P如果放贷给没有能力偿还的人，可能出现爆雷的情况。但是公募基金是不会的。

公募基金可以投资的对象受到严格的限制。比如说股票基金，只能投资上海和深圳证券交易所里交易的股票、债券等品种。不得投资非标准的金融产品。所以，爆雷的可能性很小。

但是，公募基金所投资的具体股票、债券的品种，有没有投资价值，还是需要基金经理或者投资者去判断的。

公募基金只是从产品形式上，避免了跑路和爆雷的风险。

{\heiti 基金销售平台倒闭了怎么办}

还有的朋友有疑问，我们通常投资场外基金的时候，除了传统的银行渠道，也会使用像蚂蚁财富、天天基金、蛋卷基金这样第三方基金销售平台。

那么，基金销售平台会不会没有了呢？如果没有了怎么办？

其实，这样的平台是国家认证的基金销售平台，拥有正规的基金销售牌照，在基金公司的官网上也是有授权的。

在国内，基金销售平台必须要有基金销售牌照。而基金销售牌照，国内审批也是很严格的。市面上，一个基金销售牌照的门槛也到了1亿元以上。本身拥有基金销售牌照，就已经是一件不容易的事情。

像我们熟悉的蚂蚁财富、天天基金、蛋卷基金等，都是有销售牌照的平台。通常大的银行、证券公司，也有销售牌照。

虽然我们是通过这些正规平台来进行投资，但是这些拥有基金销售牌照的平台，并不会持有我们的资金，它们只是负责销售基金。

我们通过平台来投资基金，我们的资金其实最后是到了基金公司对应的托管账户中，平台本身与我们的资金是分隔的。

因此，即使平台没有了，我们也可以在基金公司那里找到我们之前买入的基金份额。不用担心。

实际上，销售平台本身通常是轻资产运作的，也没有太大的开销，所以绝大多数销售平台可以很稳定地存在下去。

{\heiti 基金公司倒闭了怎么办}

还有一点，我们经常能听到某某基金清盘，但如果基金公司直接倒闭了，怎么办呢？

基金清盘不是说基金公司把投资者的钱吞了，而是选择某一天把基金持有者的基金份额强制赎回。

投资者可以按照某个净值来收回自己的资金，然后再去做其他的投资。

如果基金公司倒闭呢？在倒闭之前，其旗下的基金也会提前清盘了。

同样也是按照某一个净值强制赎回基金份额，投资者拿回资金后，仍然可以选择其他基金继续投资。

同样地，基金公司本身的运作成本也不高，大多数的公募基金公司也可以稳定地存在下去。

单只公募基金，如果规模在2亿元以上，对基金公司来说可以赢利，也可以长期存在下去。所以我们通常投资指数基金，也是选择规模在2亿元以上的品种。

总体来说，公募基金产品，是国内目前监管非常全面的金融产品。

证监会、银监会\textsuperscript{ ［1］ }、基金业协会、销售平台、基金公司、托管行，六方共同监管，保证公募基金产品从产品形式上风险最低。

我们投资的货币基金、债券基金、股票基金（包括指数基金），都是属于公募基金的范畴，它们在产品形式上都是没有问题的。

当然，虽然产品形式没问题，但是具体投资品种的价值，还是需要投资者自己判断。如果买的时候追高买入，还是有可能亏损的。所以公募基金投资的风险，在于投资者没有认识到投资品种的价值，容易高估的时候买入。这时收益是不理想的。

不过投资公募基金的时候，投资者不用担心产品跑路、爆雷等不必要的风险。这也是公募基金能成为全世界投资者的主要投资品种的原因。

［1］ 2018年，与保监会合并为中国银行保险监督管理委员会，简称银保监会。

\subsection{秘诀二：挑选费用较低的基金}
对指数基金收益影响比较大的，是基金的费用。

像我们在前面介绍过基金的管理费、托管费，也介绍过基金的销售服务费。

指数基金的费用主要是：申购费、赎回费、管理费、托管费。

{\heiti 透明的费率}

指数基金的申购费率，通常是1\%～1.5\%。如果直接通过银行购买，通常就是1\%～1.5\%的费率，如果通过蚂蚁财富、天天基金、蛋卷基金等平台申购，申购费一般是打1折的，也就是费率平均在0.1\%～0.15\%。这个费用其实比较少，对我们的投资影响不大。

赎回费，目前如果持有基金在7天之内进行赎回的话，会有一个赎回费，费率为1.5\%。超过7天之后，赎回费率通常是0.5\%左右。我们投资指数基金是一个长期的过程，很多指数基金在持有1～2年之后就会免除赎回费了。

管理费，也是透明的，基金报告里就有。大多数指数基金管理费率在0.5\%～0.8\%，未来还会逐渐下降，最后会在0.1\%～0.2\%。因为指数基金同质化比较严重，如果某只指数基金的管理费率更低，那么它相比同类基金就更有吸引力。所以竞争越激烈，指数基金的管理费率就会越低。

托管费，费率一般在0.1\%～0.2\%。要注意的是，所有的基金都有管理费和托管费，不管是场内基金还是场外基金。我们看到的基金净值，已经是扣除了管理费和托管费之后的。

申购费、赎回费是交易的时候收取。管理费和托管费是持有的时候，每年从基金净值中扣除。

{\heiti 不透明的费率}

不透明的费率有：分红税、印花税和交易佣金。

投资者收到基金分红没有分红税，投资者买卖基金也没有印花税。

不过基金本身会投资股票。当基金本身收到成份股的分红时，是需要交分红税的；当基金买卖成份股时，也需要交印花税、交易佣金。

这几项费用免不了，是国家或者券商收取的。

好消息是，指数基金的调仓频率是固定的，每年只能调整1～2次仓位，所以换手率相对低一些。

关于交易佣金，在基金的年报中也会有介绍。这也算是投资指数基金所面临的费率。大多数指数基金的交易佣金，一年下来大约在0.1\%～0.4\%。而主动基金，很多在1\%以上，因为主动基金的平均换手率远高于指数基金。

这部分费用，不同的基金是不同的。从整体上来看，指数基金在所有股票基金中，这部分不透明的费率是最低的。

{\heiti 各类费率的影响}

以上费用里面，对我们来说，最多的是管理费（0.5\%～0.8\%/年），其次是交易佣金（0.1\%～0.4\%/年），再次是申购赎回费（0.1\%～0.15\%/次）。其中管理费、交易佣金每年都会有。

并且交易佣金是必然会有的，这部分费用不受基金公司控制。这个费用的产生，主要是由于散户申购、赎回了基金，那么基金公司就必须要去买入、卖出对应的股票份额。

通常来说，如果频繁交易，这部分费用会很高。

只能说慢慢地让国内投资者养成长期投资的习惯，不要频繁买卖，这部分费用才会降低下来。

而管理费，随着指数基金竞争的激烈，其费率会逐渐降低。现在也有很多基金公司，愿意推出一些管理费率比较低的品种。我们应该尽力寻找费率低的优质指数基金来投资。

总体来说，指数基金的费率，是所有金融产品中最低的。只要不频繁做短线交易，这就是国内费率最低的金融产品了。其他类型的金融产品，费率远远比指数基金要高。

回到沪深300指数基金上，有的沪深300指数基金，管理费率一年可能是0.8\%，而目前管理费率最低的沪深300指数基金，管理费率是0.2\%。一年差距就有0.6\%。

如果算上托管费等其他费用，那就更高了。聪明的投资者自然会想到，如果挑选费率低的沪深300指数基金，就可以获得更好的收益了。

{\heiti 基金的费率如何查询}

在基金公司的官网，我们可以找到基金产品的页面，上面会有基金的招募书和定期报告。里面都会介绍这只基金的管理费率、托管费率等是多少。

本书介绍指数基金的章节，也有汇总国内主流指数基金的管理费率等，可供查询。

\subsection{秘诀三：选择追踪误差较小的基金}
还有一个影响指数基金收益的因素，就是追踪误差。

追踪误差是一个经常被忽略的问题。这个追踪误差，其实是指数基金特有的一个衡量指标。

我们知道指数基金追踪指数、复制指数，那指数基金的收益自然也跟指数相似。

但是有没有可能，基金公司没有追踪好指数，导致指数基金产生了比较大的误差呢？也是有可能的。

例如说指数上涨了30\%，但是指数基金只上涨了20\%，那就说明指数基金跑输了指数。

如果我们判断对了指数的投资价值，结果到头来，指数基金却产生了追踪误差，跑输了指数，那岂不是很不幸？

{\heiti 追踪误差怎么看}

这个追踪误差怎么看呢？

追踪误差在基金公告里就有写，这是每只指数基金需要定期披露的数据。

这个误差，主要是看两个方面。

第一，基金份额净值增长率与业绩比较基准收益率的差值。

基金份额净值增长率，就是这段时间里，基金净值上涨了多少。业绩比较基准收益率，就是这段时间里，指数上涨了多少。它们之间的差值，可以衡量收益上的差距。

如果基金净值增长率比业绩比较基准的收益率高，差值就是正的，就说明基金的收益比指数的默认收益好。起码，指数基金没跑输指数。

第二，基金份额净值增长率标准差与业绩比较基准收益率标准差的差值。

这个指标不常用到。标准差指的是衡量基金净值追踪的拟合程度，如果标准差的差距小，就说明基金在每个时间段的追踪情况都比较贴近于指数。

总体来说，第一个是衡量指数基金的收益是不是追得上指数，第二个则是衡量指数基金的走势是不是贴得上指数。

从我们投资指数基金的角度来说，通常前者更重要，也就是基金净值是不是跑赢了指数基准。毕竟我们更在乎收益率。

当然，第二个也很重要，它可以让我们看到每个时间段指数基金的净值是不是与指数比较接近。

我们可以在基金的定期公告中看到基金的年报和季度报，年报和季度报里都是有追踪误差情况的。

也就是我们说的“基金份额净值增长率与业绩比较基准收益率的差值”和“基金份额净值增长率标准差与业绩比较基准收益率标准差的差值”。

举一个例子，在某个月里，如果基金份额净值增长率-业绩比较基准收益率=0.74\%，说明这个月指数基金跑赢了指数0.74\%。也就是说，指数基金跑赢了指数。如果基金份额净值增长率标准差-业绩比较基准收益率标准差=0.01\%，这个值越接近0\%，说明指数基金与追踪指数贴合得越好。

{\heiti 追踪误差是好是坏}

其实A股大部分指数基金都是跑赢指数的。这是因为指数不考虑分红，而指数基金可以收到分红，这会产生一些超额收益，让指数基金的净值超越指数。

就算是本身股息率比较低的指数基金，长期来看，也有很多是跑赢指数的。这是因为国内的指数基金可以打新，打新收益会增加到基金净值里，也会让指数基金跑赢指数。

那跑赢了算不算产生了追踪误差呢？

其实严格来说，指数基金跑赢了指数，也算是产生了追踪误差。

这是一种正向的误差。对买入指数基金、上涨后卖出享受收益的投资者来说，这种误差是好事。

但是如果是跑输了指数的误差，那对买入持有的投资者就不利了。

在特殊时间段，指数基金的误差会放大，例如在牛市短期大涨的时候和熊市短期大跌的时候。

牛市短期大涨，会有大量的投资者，抢着申购指数基金。这些资金到了基金公司那里，它们买入股票是需要时间的。但这时可能股票短期里仍然在大涨，这就导致牛市里部分指数基金会暂时跑输指数。这种情况基金公司也没有太好的办法。

还有一种情况，是熊市短期大跌的时候。牛市过去之后，可能市场会有一些短期大幅下跌。例如2015年6月牛市达到了最顶峰，之后几个月市场短期大跌。下跌过程中，很多上市公司临时停牌。停牌后，这些公司的股票是无法交易的，指数基金持有这些公司的股票，自然也无法卖出。这也会对指数基金的追踪产生一定误差。

不过市场不会一直上涨，股票也不会一直停盘。当这些时间段过去，指数基金的追踪又会回归正常。所以遇到短期出现追踪误差时，如果是上面两种原因，也不用过于担心。

总体来说，规模至少大于1亿元、管理费率和托管费率等越低、追踪误差控制良好的指数基金，是我们主要考虑的对象。

对规模、追踪误差的要求，主要是减少一些不必要的风险。

费率低了，长期来看可以提高0.3\%～0.5\%的年化收益率，非常有效地增加我们的收益。

挑选指数基金品种，这是一个基础的优化我们投资收益的操作。

这些并不复杂，只要知道了对应基金品种的规模、费率、追踪误差等信息，就可以做出决策。

定投指数基金还有一些其他的技巧，可以帮助我们提高收益。我们将在接下来的几个章节中，为大家一一介绍。

{\heiti 投资者笔记}

\begin{itemize}
\item {\kaishu 挑选指数基金的时候，通常有三个比较重要的因素，分别是规模、费用和追踪误差。}
\item {\kaishu 值得考虑的指数基金，通常符合这三个条件：规模至少大于1亿元，管理费率、托管费率等越低越好，追踪误差控制良好。}
\item {\kaishu 普通投资者投资公募基金时，不用担心产品跑路、爆雷等不必要的风险。}
\end{itemize}

